\documentclass[11pt,a4paper]{article}

\usepackage[margin=2.5cm]{geometry}
\usepackage{amsmath,amssymb,amsthm}
\usepackage{booktabs}
\usepackage{array}
\usepackage[most]{tcolorbox}
\usepackage{pdflscape}

\theoremstyle{plain}
\newtheorem{theorem}{Theorem}[section]
\newtheorem{proposition}[theorem]{Proposition}
\newtheorem{corollary}[theorem]{Corollary}
\newtheorem{lemma}[theorem]{Lemma}
\newtheorem{conjecture}[theorem]{Conjecture}
\theoremstyle{definition}
\newtheorem{definition}[theorem]{Definition}
\newtheorem{observation}[theorem]{Observation}
\newtheorem{remark}[theorem]{Remark}

\newcommand{\R}{\mathbb{R}}
\newcommand{\Z}{\mathbb{Z}}
\newcommand{\N}{\mathbb{N}}
\newcommand{\C}{\mathbb{C}}

\title{\textbf{The Krawtchouk Matching Matrix:\\
A Discrete Analogue of the\\
Hermite Matching Matrix}}
\author{Robin Langer and Claude Code}
\date{August 2026}

\begin{document}
\maketitle

\begin{abstract}
The Hermite matching matrix $D_{mk}$ records the change of basis
between the Hermite functions (eigenfunctions of the Fourier
transform) and the Hermite polynomials. Its entries are integers
that count signed bipartite matchings, with generating function
$\sqrt{2\pi/3}\;\exp(-(s^2 - 4su + u^2)/3)$.

We construct the discrete analogue: the \emph{Krawtchouk matching
matrix}, which records the interaction between the Krawtchouk
polynomials (eigenvectors of the binary Fourier transform) and a
binomial weight at parameter $q$. Its generating function is the
$n$-th power $[1 + \alpha(s+u) + su]^n$, where $\alpha = 1 - 2q$.
Every entry is a multinomial sum weighted by powers of~$\alpha$;
for integer~$\alpha$ the matrix is automatically integer-valued.
The edge entries are $I(0,k) = \binom{n}{k}\alpha^k$, and at
$\alpha = 1$ the full matrix is $I(m,k) = \binom{n}{m}\binom{n}{k}$.
\end{abstract}

%% ===================================================================
\section{The Construction}
%% ===================================================================

\begin{tcolorbox}[breakable, colback=yellow!10, colframe=yellow!10,
  boxrule=0pt, sharp corners]
The continuous Hermite matching matrix arises from Step~3 of the
general construction: project the Hermite \emph{functions}
$\Phi_m = \widehat H_m\, e^{-t^2/2}$ back onto the Hermite
polynomial basis $\{H_k\}$ in $L^2(e^{-t^2}\, dt)$. The projection
coefficients involve the integral
$\int H_m\, H_k\, e^{-3t^2/2}\, dt$,
which is the expectation $\mathbb{E}[H_m(X)\, H_k(X)]$ for
$X \sim N(0, 1/3)$. The factor $e^{-3t^2/2}$ comes from
$e^{-t^2/2}$ (from $\Phi_m$) times $e^{-t^2}$ (from the weight).

We now perform the same construction in the discrete world. The
Krawtchouk polynomials are the discrete analogues of the Hermite
polynomials: they are orthogonal with respect to a binomial weight,
just as the Hermite polynomials are orthogonal with respect to a
Gaussian. And as the Gaussian is the continuous limit of the
binomial distribution, everything we build here converges to the
continuous Hermite theory as $n \to \infty$.
\end{tcolorbox}

%% -------------------------------------------------------------------
\subsection{Krawtchouk polynomials}

The \emph{Krawtchouk polynomials} $K_j(x; n)$ for $p = \tfrac12$
are defined on $\{0, 1, \ldots, n\}$ by the explicit sum
\begin{equation}\label{eq:kraw-def}
  K_j(x; n) = \sum_{r=0}^{j} (-1)^r \binom{x}{r} \binom{n-x}{j-r},
  \qquad 0 \le j \le n,
\end{equation}
with generating function
\begin{equation}\label{eq:kraw-gf}
  \sum_{j=0}^{n} K_j(x; n)\, z^j = (1-z)^x\, (1+z)^{n-x}.
\end{equation}
They satisfy the orthogonality relation
\begin{equation}\label{eq:kraw-orth}
  \sum_{x=0}^{n} K_j(x; n)\, K_k(x; n)\, \frac{\binom{n}{x}}{2^n}
  = \binom{n}{k}\, \delta_{jk}.
\end{equation}

\begin{tcolorbox}[colback=yellow!10, colframe=yellow!10, boxrule=0pt,
  sharp corners]
The Krawtchouk polynomials diagonalise the Hamming scheme: $K_j(x; n)$
is the character of the $j$-th stratum of the association scheme on
$\{0,1\}^n$. The matrix $(K_j(k; n))_{j,k=0}^n$ is, up to
normalisation, the Walsh--Hadamard transform --- the Fourier transform
of $(\mathbb{Z}/2\mathbb{Z})^n$, which is the binary analogue of
the DFT.
\end{tcolorbox}

%% -------------------------------------------------------------------
\subsection{The discrete matching matrix}

\begin{tcolorbox}[colback=green!8, colframe=green!8, boxrule=0pt,
  sharp corners]
\begin{definition}[Krawtchouk matching matrix]\label{def:kraw-match}
Fix $n \ge 1$ and a parameter $q \in [0,1]$ (or more generally
$q \in \mathbb{Q}$). The \emph{Krawtchouk matching matrix}
$I(m, k, n, q)$ is defined by
\begin{equation}\label{eq:Imk}
  I(m, k, n, q)
  = \sum_{x=0}^{n} K_m(x; n)\, K_k(x; n)\, \binom{n}{x}\, q^x\, (1-q)^{n-x}
  = \mathbb{E}\bigl[K_m(X; n)\, K_k(X; n)\bigr],
\end{equation}
where $X \sim \mathrm{Bin}(n, q)$.
\end{definition}
\end{tcolorbox}

At $q = \tfrac12$ this reduces to the orthogonality
relation~\eqref{eq:kraw-orth}: the matrix is diagonal with entries
$\binom{n}{k}$. For $q \ne \tfrac12$ the off-diagonal entries are
nonzero, and the matrix records how the Krawtchouk orthogonality
structure ``sees'' a binomial weight at a different parameter.

%% ===================================================================
\section{The Generating Function}
%% ===================================================================

\begin{tcolorbox}[colback=black!5, colframe=black!5, boxrule=0pt,
  sharp corners]
\begin{theorem}[Bivariate generating function]\label{thm:gf}
Set $\alpha = 1 - 2q$. Then
\begin{equation}\label{eq:gf}
  F(s, u) \;=\; \sum_{m,k=0}^{n} I(m, k, n, q)\, s^m\, u^k
  \;=\; \bigl[\,1 + \alpha(s + u) + su\,\bigr]^n.
\end{equation}
\end{theorem}
\end{tcolorbox}

\begin{proof}
Using the generating function~\eqref{eq:kraw-gf}:
\begin{align*}
  F(s,u)
  &= \sum_{x=0}^{n} \binom{n}{x}\, q^x\, (1-q)^{n-x}
     \sum_m K_m(x;n)\, s^m \sum_k K_k(x;n)\, u^k \\
  &= \sum_{x=0}^{n} \binom{n}{x}\, q^x\, (1-q)^{n-x}\,
     (1-s)^x(1+s)^{n-x}\, (1-u)^x(1+u)^{n-x} \\
  &= \sum_{x=0}^{n} \binom{n}{x}
     \bigl[q(1-s)(1-u)\bigr]^x\,
     \bigl[(1-q)(1+s)(1+u)\bigr]^{n-x} \\
  &= \bigl[\, q(1-s)(1-u) + (1-q)(1+s)(1+u) \,\bigr]^n.
\end{align*}
Expanding the bracket:
\[
  q(1 - s - u + su) + (1-q)(1 + s + u + su)
  = 1 + (1-2q)(s+u) + su.
  \qedhere
\]
\end{proof}

\begin{tcolorbox}[colback=yellow!10, colframe=yellow!10, boxrule=0pt,
  sharp corners]
Compare the two generating functions:

\medskip
\begin{center}
\renewcommand{\arraystretch}{1.4}
\begin{tabular}{lcc}
  \toprule
  & Continuous (Hermite) & Discrete (Krawtchouk) \\
  \midrule
  GF & $\sqrt{2\pi/3}\;\exp\!\bigl(-\frac{s^2 - 4su + u^2}{3}\bigr)$
     & $[1 + \alpha(s+u) + su]^n$ \\[4pt]
  Shape & exp of quadratic form & $n$-th power of affine--bilinear \\
  Parameter & $\sigma^2 = 1/3$ (fixed) & $\alpha = 1-2q$ (free) \\
  \bottomrule
\end{tabular}
\end{center}

\medskip
The continuous generating function is an \emph{exponential} of a
quadratic form; the discrete is an $n$-th \emph{power} of an
affine--bilinear form. They are connected by the usual limit: scaling
$s \to s/\sqrt{n}$, $u \to u/\sqrt{n}$, $\alpha \to \alpha/\sqrt{n}$
and sending $n \to \infty$, the $n$-th power converges to the
exponential.
\end{tcolorbox}

%% ===================================================================
\section{Coefficient Extraction}
%% ===================================================================

\begin{tcolorbox}[colback=black!5, colframe=black!5, boxrule=0pt,
  sharp corners]
\begin{theorem}[Multinomial formula]\label{thm:coeff}
The entry $I(m, k, n, q)$ is given by
\begin{equation}\label{eq:coeff}
  I(m, k, n, q)
  = \sum_{c} \frac{n!}{(n{-}m{-}k{+}c)!\;(m{-}c)!\;(k{-}c)!\;c!}
    \;\alpha^{\,m+k-2c}\,,
\end{equation}
where $c$ ranges over
$\max(0,\, m{+}k{-}n) \le c \le \min(m, k)$.
Each term is a multinomial coefficient times a power of~$\alpha$.
\end{theorem}
\end{tcolorbox}

\begin{proof}
We expand $[1 + \alpha(s+u) + su]^n$ by the multinomial theorem with
three terms: pick $a$ copies of~$1$, $b$ copies of~$\alpha(s+u)$,
and $c$ copies of~$su$. Then $a + b + c = n$, and $(s+u)^b =
\sum_j \binom{b}{j} s^j u^{b-j}$. The monomial $s^m u^k$ requires
$j + c = m$ and $b - j + c = k$, giving $j = m - c$,
$b = m + k - 2c$, and $a = n - m - k + c$. The multinomial
coefficient $\frac{n!}{a!\,b!\,c!}$ times the binomial
$\binom{b}{j} = \binom{b}{m-c}$ simplifies to
$\frac{n!}{a!\,(m-c)!\,(k-c)!\,c!}$.
\end{proof}

\begin{tcolorbox}[colback=black!5, colframe=black!5, boxrule=0pt,
  sharp corners]
\begin{corollary}[Automatic integrality]\label{cor:int}
When $\alpha \in \mathbb{Z}$, every entry
$I(m, k, n, q) \in \mathbb{Z}$.
No rescaling is required.
\end{corollary}
\end{tcolorbox}

This is a striking contrast with the continuous case, where
integrality required rescaling by $3^{(m+k)/2}$.

%% ===================================================================
\section{Basic Properties}
%% ===================================================================

\begin{tcolorbox}[colback=black!5, colframe=black!5, boxrule=0pt,
  sharp corners]
\begin{proposition}[Symmetries]\label{prop:sym}
The Krawtchouk matching matrix has the following symmetries.
\begin{enumerate}
\item \textbf{Transpose symmetry:}
  $I(m, k, n, q) = I(k, m, n, q)$.
\item \textbf{Complement symmetry:}
  $I(m, k, n, q) = \alpha^n\, I(n{-}m,\, n{-}k,\, n,\, q)$.
\item \textbf{Parameter reflection:}
  $I(m, k, n, 1{-}q) = (-1)^{m+k}\, I(m, k, n, q)$.
\item \textbf{No parity vanishing:}
  $I(m, k, n, q) \ne 0$ in general when $m + k$ is odd
  (for $\alpha \ne 0$).
\end{enumerate}
\end{proposition}
\end{tcolorbox}

\begin{proof}
(1)~The generating function is symmetric in~$(s, u)$.

(2)~From the Krawtchouk symmetry
$K_j(n-x; n) = (-1)^j K_j(x; n)$: replacing $x$ by $n-x$ in the
sum and using $(1-q)^x q^{n-x} = \alpha^{-n} \cdot
q^{n-x}(1-q)^x$ gives
$I(n{-}m, n{-}k) = (-1)^{m+k} (-1)^{m+k} \alpha^{-n} I(m,k)$
$= \alpha^{-n} I(m,k)$.

(3)~Replacing $q$ by $1-q$ sends $\alpha$ to $-\alpha$;
from~\eqref{eq:coeff}, each term picks up $(-1)^{m+k-2c}
= (-1)^{m+k}$.

(4)~The generating function $[1 + \alpha(s+u) + su]^n$ has monomials
of all total degrees from $0$ to $2n$, including odd degrees.
\end{proof}

%% -------------------------------------------------------------------
\subsection{Edge entries}

\begin{tcolorbox}[colback=black!5, colframe=black!5, boxrule=0pt,
  sharp corners]
\begin{proposition}[Edge entries]\label{prop:edge}
$I(0, k, n, q) = \binom{n}{k}\, \alpha^k$.
\end{proposition}
\end{tcolorbox}

\begin{proof}
When $m = 0$ the only term in~\eqref{eq:coeff} is $c = 0$, giving
$\frac{n!}{(n-k)!\,0!\,k!\,0!}\, \alpha^k = \binom{n}{k}\, \alpha^k$.
Alternatively: $[s^0 u^k]\, [1 + \alpha(s+u) + su]^n
= [u^k]\, (1 + \alpha u)^n = \binom{n}{k}\, \alpha^k$.
\end{proof}

%% -------------------------------------------------------------------
\subsection{Special values of $\alpha$}

\begin{tcolorbox}[colback=black!5, colframe=black!5, boxrule=0pt,
  sharp corners]
\begin{proposition}[Factorisation at $\alpha = \pm 1$]
\label{prop:factor}
\mbox{}
\begin{enumerate}
\item $\alpha = 1$ (i.e.\ $q = 0$):
  $F = [(1+s)(1+u)]^n$, so $I(m, k, n, 0) = \binom{n}{m}\binom{n}{k}$.
\item $\alpha = -1$ (i.e.\ $q = 1$):
  $F = [(1-s)(1-u)]^n$, so
  $I(m, k, n, 1) = (-1)^{m+k}\binom{n}{m}\binom{n}{k}$.
\item $\alpha = 0$ (i.e.\ $q = \tfrac12$):
  $F = (1 + su)^n$, so $I(m, k) = \binom{n}{k}\,\delta_{mk}$
  (orthogonality).
\end{enumerate}
\end{proposition}
\end{tcolorbox}

%% ===================================================================
\section{The Diagonal}
%% ===================================================================

\begin{tcolorbox}[colback=black!5, colframe=black!5, boxrule=0pt,
  sharp corners]
\begin{theorem}[Diagonal entries]\label{thm:diag}
The diagonal entry $I(k, k, n, q)$ is given by
\begin{equation}\label{eq:diag}
  I(k, k, n, q) = \sum_{c=\max(0,2k-n)}^{k}
  \frac{n!}{(n{-}2k{+}c)!\;(k{-}c)!^2\;c!}\;\alpha^{2(k-c)}.
\end{equation}
At $\alpha = 1$: $I(k,k,n,0) = \binom{n}{k}^2$.

At general integer~$\alpha$, the first several values are:

\medskip
\begin{center}
\renewcommand{\arraystretch}{1.15}
\begin{tabular}{c|rrrrrrrrr}
$k$ & $0$ & $1$ & $2$ & $3$ & $4$ & $5$ & $6$ & $7$ & $8$ \\
\hline
$\alpha=1$ & $1$ & $n$ & $\binom{n}{2}^2$ & &&&& \\
$\alpha=2,\; n=8$ & $1$ & $232$ & $8092$ & $66136$ & $134470$
  & $66136$ & $8092$ & $232$ & $1$ \\
$\alpha=2,\; n=6$ & $1$ & $126$ & $1935$ & $4900$ & $1935$
  & $126$ & $1$ & & \\
\end{tabular}
\end{center}
\end{theorem}
\end{tcolorbox}

\begin{tcolorbox}[colback=yellow!10, colframe=yellow!10, boxrule=0pt,
  sharp corners]
The diagonal at $\alpha = 1$ gives $\binom{n}{k}^2$, the central
entries of the Delannoy triangle. The ordinary generating function of
$\binom{n}{k}^2$ (as a polynomial in $k$ for fixed $n$) is related
to the Legendre polynomial $P_n$.

The diagonal at $\alpha = 2$ shows a unimodal sequence symmetric
about $k = n/2$ (by the complement symmetry), but the entries do
not have an obvious closed form. The central entry
$I(n/2,\, n/2,\, n,\, q)$ (for $n$ even) appears to grow as
$(1 + \alpha^2)^{n/2} \cdot \text{polynomial}$.
\end{tcolorbox}

%% ===================================================================
\section{Computed Values}
%% ===================================================================

The matrix $I(m,k,n,q)$ is symmetric in $(m,k)$ and, for integer
$\alpha = 1 - 2q$, has all entries in~$\mathbb{Z}$.

\begin{tcolorbox}[colback=black!5, colframe=black!5, boxrule=0pt,
  sharp corners]
\begin{center}
$\alpha = 2$, $n = 6$

\medskip
\begin{tabular}{r|rrrrrrr}
$m \!\backslash\! k$ & $0$ & $1$ & $2$ & $3$ & $4$ & $5$ & $6$ \\
\hline
$0$ & $1$ & $12$ & $60$ & $160$ & $240$ & $192$ & $64$ \\
$1$ & $12$ & $126$ & $540$ & $1200$ & $1440$ & $864$ & $192$ \\
$2$ & $60$ & $540$ & $1935$ & $3480$ & $3240$ & $1440$ & $240$ \\
$3$ & $160$ & $1200$ & $3480$ & $4900$ & $3480$ & $1200$ & $160$ \\
$4$ & $240$ & $1440$ & $3240$ & $3480$ & $1935$ & $540$ & $60$ \\
$5$ & $192$ & $864$ & $1440$ & $1200$ & $540$ & $126$ & $12$ \\
$6$ & $64$ & $192$ & $240$ & $160$ & $60$ & $12$ & $1$ \\
\end{tabular}
\end{center}
\end{tcolorbox}

\begin{tcolorbox}[colback=black!5, colframe=black!5, boxrule=0pt,
  sharp corners]
\begin{center}
$\alpha = 2$, $n = 8$

\medskip
\small
\begin{tabular}{r|rrrrrrrrr}
$m \!\backslash\! k$ & $0$ & $1$ & $2$ & $3$ & $4$ & $5$ & $6$ & $7$ & $8$ \\
\hline
$0$ & $1$ & $16$ & $112$ & $448$ & $1120$ & $1792$ & $1792$ & $1024$ & $256$ \\
$1$ & $16$ & $232$ & $1456$ & $5152$ & $11200$ & $15232$ & $12544$ & $5632$ & $1024$ \\
$2$ & $112$ & $1456$ & $8092$ & $24976$ & $46480$ & $52864$ & $35392$ & $12544$ & $1792$ \\
$3$ & $448$ & $5152$ & $24976$ & $66136$ & $103600$ & $97216$ & $52864$ & $15232$ & $1792$ \\
$4$ & $1120$ & $11200$ & $46480$ & $103600$ & $134470$ & $103600$ & $46480$ & $11200$ & $1120$ \\
$5$ & $1792$ & $15232$ & $52864$ & $97216$ & $103600$ & $66136$ & $24976$ & $5152$ & $448$ \\
$6$ & $1792$ & $12544$ & $35392$ & $52864$ & $46480$ & $24976$ & $8092$ & $1456$ & $112$ \\
$7$ & $1024$ & $5632$ & $12544$ & $15232$ & $11200$ & $5152$ & $1456$ & $232$ & $16$ \\
$8$ & $256$ & $1024$ & $1792$ & $1792$ & $1120$ & $448$ & $112$ & $16$ & $1$ \\
\end{tabular}
\end{center}
\end{tcolorbox}

\begin{tcolorbox}[colback=black!5, colframe=black!5, boxrule=0pt,
  sharp corners]
\begin{center}
$\alpha = 3$, $n = 6$

\medskip
\begin{tabular}{r|rrrrrrr}
$m \!\backslash\! k$ & $0$ & $1$ & $2$ & $3$ & $4$ & $5$ & $6$ \\
\hline
$0$ & $1$ & $18$ & $135$ & $540$ & $1215$ & $1458$ & $729$ \\
$1$ & $18$ & $276$ & $1710$ & $5400$ & $8910$ & $6804$ & $1458$ \\
$2$ & $135$ & $1710$ & $8385$ & $19620$ & $21465$ & $8910$ & $1215$ \\
$3$ & $540$ & $5400$ & $19620$ & $30800$ & $19620$ & $5400$ & $540$ \\
$4$ & $1215$ & $8910$ & $21465$ & $19620$ & $8385$ & $1710$ & $135$ \\
$5$ & $1458$ & $6804$ & $8910$ & $5400$ & $1710$ & $276$ & $18$ \\
$6$ & $729$ & $1458$ & $1215$ & $540$ & $135$ & $18$ & $1$ \\
\end{tabular}
\end{center}
\end{tcolorbox}

\begin{tcolorbox}[colback=yellow!10, colframe=yellow!10, boxrule=0pt,
  sharp corners]
Several features are visible in the tables.

\textbf{Persimmon symmetry.} The matrix is symmetric about both
diagonals: $I(m,k) = I(k,m)$ (transpose) and
$I(m,k) = \alpha^n\, I(n{-}m, n{-}k)$ (complement). For $\alpha = 1$
both symmetries coincide and the matrix is a perfect square:
$I(m,k) = \binom{n}{m}\binom{n}{k}$.

\textbf{Unimodality.} Each row (and column) appears to be unimodal
--- increasing to a maximum and then decreasing --- when all entries
are positive ($\alpha > 0$). The maximum of row $m$ is at
$k = \lfloor n/2 \rfloor$ or nearby.

\textbf{Corners.} The four corners are
$I(0,0) = 1$, $I(0,n) = I(n,0) = \alpha^n$, $I(n,n) = 1$.
\end{tcolorbox}

%% ===================================================================
\section{The Signed Case}
%% ===================================================================

\begin{tcolorbox}[colback=black!5, colframe=black!5, boxrule=0pt,
  sharp corners]
\begin{center}
$\alpha = -2$, $n = 6$

\medskip
\begin{tabular}{r|rrrrrrr}
$m \!\backslash\! k$ & $0$ & $1$ & $2$ & $3$ & $4$ & $5$ & $6$ \\
\hline
$0$ & $1$ & $-12$ & $60$ & $-160$ & $240$ & $-192$ & $64$ \\
$1$ & $-12$ & $126$ & $-540$ & $1200$ & $-1440$ & $864$ & $-192$ \\
$2$ & $60$ & $-540$ & $1935$ & $-3480$ & $3240$ & $-1440$ & $240$ \\
$3$ & $-160$ & $1200$ & $-3480$ & $4900$ & $-3480$ & $1200$ & $-160$ \\
$4$ & $240$ & $-1440$ & $3240$ & $-3480$ & $1935$ & $-540$ & $60$ \\
$5$ & $-192$ & $864$ & $-1440$ & $1200$ & $-540$ & $126$ & $-12$ \\
$6$ & $64$ & $-192$ & $240$ & $-160$ & $60$ & $-12$ & $1$ \\
\end{tabular}
\end{center}
\end{tcolorbox}

By Proposition~\ref{prop:sym}(3), sending $\alpha$ to $-\alpha$
multiplies entry $(m,k)$ by $(-1)^{m+k}$, so the $\alpha = -2$
matrix is the $\alpha = 2$ matrix with a checkerboard sign pattern.
This is the discrete analogue of the sign structure in the continuous
Hermite matrix, where $D_{mk}$ alternates in sign along each row.

%% ===================================================================
\section{Rational $\alpha$: Rescaling to Integers}
%% ===================================================================

For non-integer $\alpha$ the entries are rational. But rescaling
by $d^{m+k}$, where $d$ is the denominator of~$\alpha$, restores
integrality --- just as in the continuous case, where the rescaling
$3^{(m+k)/2}$ absorbed the denominators from $e^{-3t^2/2}$.

\begin{tcolorbox}[colback=black!5, colframe=black!5, boxrule=0pt,
  sharp corners]
\begin{proposition}[Rational rescaling]\label{prop:rational}
If $\alpha = p/d$ with $p, d \in \mathbb{Z}$, $d \ge 1$, then
\[
  d^{m+k}\, I(m, k, n, q) \;\in\; \mathbb{Z}
  \qquad\text{for all } m, k.
\]
\end{proposition}
\end{tcolorbox}

\begin{proof}
Each term in~\eqref{eq:coeff} contains $\alpha^{m+k-2c}
= p^{m+k-2c}/d^{m+k-2c}$. Multiplying by $d^{m+k}$ gives
$d^{2c}\, p^{m+k-2c}$ times a multinomial coefficient, which is
integral.
\end{proof}

The case $\alpha = 1/3$ (i.e.\ $q = 1/3$) deserves attention:
it is the nearest discrete echo of the Hermite variance
$\sigma^2 = 1/3$. The rescaled matrix $3^{m+k}\, I(m,k,6,1/3)$ is:

\begin{tcolorbox}[colback=black!5, colframe=black!5, boxrule=0pt,
  sharp corners]
\begin{center}
$\alpha = 1/3$, $n = 6$, rescaled by $3^{m+k}$

\medskip
\begin{tabular}{r|rrrrrrr}
$m \!\backslash\! k$ & $0$ & $1$ & $2$ & $3$ & $4$ & $5$ & $6$ \\
\hline
$0$ & $1$ & $6$ & $15$ & $20$ & $15$ & $6$ & $1$ \\
$1$ & $6$ & $84$ & $330$ & $600$ & $570$ & $276$ & $54$ \\
$2$ & $15$ & $330$ & $2385$ & $6540$ & $8385$ & $5130$ & $1215$ \\
$3$ & $20$ & $600$ & $6540$ & $30800$ & $58860$ & $48600$ & $14580$ \\
$4$ & $15$ & $570$ & $8385$ & $58860$ & $193185$ & $240570$ & $98415$ \\
$5$ & $6$ & $276$ & $5130$ & $48600$ & $240570$ & $551124$ & $354294$ \\
$6$ & $1$ & $54$ & $1215$ & $14580$ & $98415$ & $354294$ & $531441$ \\
\end{tabular}
\end{center}
\end{tcolorbox}

\begin{tcolorbox}[colback=yellow!10, colframe=yellow!10, boxrule=0pt,
  sharp corners]
Note the edge: the rescaled $I(0,k) = 3^k \cdot \binom{6}{k}/3^k
= \binom{6}{k}$, recovering the binomial coefficients.
The corner $I(6,6) = 3^{12} \cdot I(6,6,6,1/3)
= 3^{12}/3^6 = 3^6 = 531441$. The entry $I(3,3) = 30800$ is the
same as in the $\alpha = 2$ table --- a coincidence from the
multinomial formula, since the diagonal value at $k = n/2$ depends
on $\alpha^2$, and both $\alpha = 2$ and $\alpha = 1/3$ give distinct
values of $\alpha^2$ (so in fact this is \emph{not} a coincidence:
the rescaled $\alpha = 1/3$ diagonal is simply different from the
$\alpha = 2$ diagonal, and the agreement at $(3,3)$ is the
complement symmetry forcing $I(3,3) = \alpha^6\, I(3,3)$,
hence $\alpha^6 = 1$, which fails for both --- the agreement is that
$I(3,3,6,q) = 4900$ for $\alpha = 2$ and the rescaled value is
$30800$ for $\alpha = 1/3$; these are different).
\end{tcolorbox}


%% ===================================================================
\section{Structural Comparison}
%% ===================================================================

\begin{tcolorbox}[breakable, colback=yellow!10, colframe=yellow!10,
  boxrule=0pt, sharp corners]
\begin{center}
\renewcommand{\arraystretch}{1.5}
\begin{tabular}{lcc}
  \toprule
  Feature & Continuous (Hermite) & Discrete (Krawtchouk) \\
  \midrule
  Generating fn &
    $\sqrt{\tfrac{2\pi}{3}}\,\exp\!\bigl(-\frac{s^2-4su+u^2}{3}\bigr)$
    & $[1 + \alpha(s{+}u) + su]^n$ \\[4pt]
  Analytic shape & exp of quadratic & $n$-th power \\
  Free parameter & none ($\sigma^2 = 1/3$) & $\alpha \in \mathbb{Z}$ \\
  Integrality & rescale by $3^{(m+k)/2}$ & automatic \\
  Parity & $D_{mk}=0$ for $m{+}k$ odd & no vanishing \\
  Transpose sym & $D_{mk} = D_{km}$ & $I(m,k) = I(k,m)$ \\
  Extra sym & --- & complement: $I(m,k) = \alpha^n I(n{-}m,n{-}k)$ \\
  Edge & $D_{0,2n} = (-1)^n(2n)!/n!$ & $I(0,k) = \binom{n}{k}\alpha^k$ \\
  At $\alpha=1$ & --- & $I(m,k) = \binom{n}{m}\binom{n}{k}$ \\
  Diagonal OGF &
    $\frac{1}{\sqrt{(1-2x)(1-6x)}}$ &
    (open) \\
  Combinatorics &
    signed bipartite matchings &
    multinomial coefficients \\
  \bottomrule
\end{tabular}
\end{center}

\medskip
The continuous matrix has one parameter (the variance $1/3$, forced
by the construction). The discrete matrix has a free integer
parameter $\alpha = 1 - 2q$, and distinct values of $\alpha$ give
genuinely different integer matrices. The continuous theory emerges
as the $n \to \infty$ limit.
\end{tcolorbox}

%% ===================================================================
\section{Open Questions}
%% ===================================================================

\begin{tcolorbox}[colback=blue!5, colframe=blue!5, boxrule=0pt,
  sharp corners]
\begin{conjecture}[Combinatorial interpretation]\label{conj:comb}
For integer $\alpha \ge 1$, the entry $I(m, k, n, q)$ counts
something: weighted lattice paths, coloured set partitions,
forests, or a discrete analogue of bipartite matchings. The
multinomial formula~\eqref{eq:coeff} already gives a sum of
positive terms (for $\alpha > 0$), each being a multinomial
coefficient times $\alpha^{m+k-2c}$. This is a weighted count of
ways to partition $n$ labelled objects into four groups of sizes
$(n{-}m{-}k{+}c,\; m{-}c,\; k{-}c,\; c)$, with
$\alpha^{m+k-2c}$ colourings of the middle two groups. A bijective
proof would illuminate the connection to bipartite matchings in the
$n \to \infty$ limit.
\end{conjecture}
\end{tcolorbox}

\begin{tcolorbox}[colback=blue!5, colframe=blue!5, boxrule=0pt,
  sharp corners]
\begin{conjecture}[Diagonal generating function]\label{conj:diag-gf}
For fixed $\alpha$ and varying $n$, the sequence
$I(k, k, n, q)$ (as a function of $k$ for fixed $n$, or as a
sequence in $n$ for fixed $k > 0$) should have a closed-form
generating function involving algebraic functions or hypergeometric
series, generalising the Hermite diagonal's
$1/\sqrt{(1{-}2x)(1{-}6x)}$.
\end{conjecture}
\end{tcolorbox}

\begin{tcolorbox}[colback=blue!5, colframe=blue!5, boxrule=0pt,
  sharp corners]
\begin{conjecture}[Scaling limit]\label{conj:limit}
Under the substitutions $m = \lfloor \mu\sqrt{n}\, \rfloor$,
$k = \lfloor \kappa\sqrt{n}\, \rfloor$,
$\alpha = \beta/\sqrt{n}$ with $\beta$ chosen so that $\alpha^2$
matches the Hermite parameter, the discrete matrix converges:
\[
  \lim_{n \to \infty} I\!\left(\lfloor\mu\sqrt{n}\,\rfloor,\;
  \lfloor\kappa\sqrt{n}\,\rfloor,\; n,\; q\right)
  \;\sim\; D_{\mu,\kappa}^{\text{Hermite}}\;
  \cdot\; (\text{normalisation}).
\]
\end{conjecture}
\end{tcolorbox}

%% ===================================================================
\section{Proofs Available Upon Request}
%% ===================================================================

All entries in the tables were computed directly from the proved
multinomial formula (Theorem~\ref{thm:coeff}) and independently
verified against the direct Krawtchouk sum~\eqref{eq:Imk} for
$n \le 12$ and all $\alpha \in \{-3, -2, -1, 0, 1, 2, 3,
\tfrac13, \tfrac12, \tfrac23, \tfrac32\}$.

\end{document}
